\newgeometry{margin=.5cm}
\pagestyle{empty}
\raggedcolumns
\footnotesize
\newcommand{\cheatSheetColor}{RoyalBlue}

\newtcolorbox{cheatSheetBox}[1][]{
enhanced jigsaw,
% breakable,
colback=\cheatSheetColor!10,
colframe=\cheatSheetColor,
colbacktitle=\cheatSheetColor!30,
fonttitle=\bfseries,
parbox=false,
attach boxed title to top text left={yshift=-3mm,yshifttext=-3mm},
coltitle=black,
#1}

\lstset{style=python, basicstyle=\ttfamily\footnotesize}

\begin{landscape}
\footnotesize
\begin{multicols*}{4}
\textcolor{\cheatSheetColor}{\begin{center}
\section{Cheatsheet}
\end{center}\hrule}

\textcolor{\cheatSheetColor}{\subsection*{Orientierung}}

\begin{cheatSheetBox}[title=Funktionsarten]
\begin{itemize}
    \item \textbf{Priorisierung:} Elemente werden nach Relevanz oder Wichtigkeit geordnet.
    \item \textbf{Klassifikation:} Ein Datenpunkt wird einer Kategorie zugeordnet.
    \item \textbf{Assoziation:} Häufig gemeinsam auftretende Muster werden gefunden.
    \item \textbf{Filterung:} Inhalte werden zugelassen, sortiert oder ausgeblendet.
\end{itemize}
\end{cheatSheetBox}

\begin{cheatSheetBox}[title=Zentrale Formeln]
\[
\text{Score}(x)=\sum_{i=1}^{n} w_i z_i
\qquad
PR(p)=\frac{1-d}{N}+d\sum_{q\to p}\frac{PR(q)}{\text{outdeg}(q)}
\qquad
\text{Gini}(t)=1-\sum_k p_k^2
\]
Gewichte, Linkstruktur und Klassenanteile bestimmen die jeweilige Priorisierung oder Trennung.
\end{cheatSheetBox}

\columnbreak\textcolor{\cheatSheetColor}{\subsection*{Klassifikation}}

\begin{cheatSheetBox}[title=Entscheidungsbaum]
Ein Baum stellt nacheinander Fragen an die Daten und teilt sie durch Schwellenwerte auf.
Wichtige Begriffe sind Knoten, Kante, Blatt und Unreinheit.

\[
\text{Gini}(t)=1-\sum_k p_k^2
\]

Je reiner ein Knoten, desto kleiner ist der Gini-Wert.
\end{cheatSheetBox}

\begin{cheatSheetBox}[title=Random Forest]
Random Forest kombiniert viele zufällige Bäume und entscheidet per Mehrheitsentscheid.
Wichtige Einstellungen sind die Anzahl der Bäume, die Tiefe und die Blattgrösse.
Er ist robuster als ein einzelner Baum, aber weniger direkt erklärbar.
\end{cheatSheetBox}

\begin{cheatSheetBox}[title=KNN]
KNN vergleicht einen neuen Punkt mit seinen nächsten Nachbarn und übernimmt deren Mehrheitsklasse.
Wichtig ist die Normalisierung, weil Distanzen von der Skala der Features abhängen.

\[
d(a,b)=\sqrt{\sum_{i=1}^{n}(a_i-b_i)^2}
\]

Kleines $k$ reagiert empfindlich, grosses $k$ glättet die Entscheidung.
\end{cheatSheetBox}

\begin{cheatSheetBox}[title=SVM als Bonus]
Eine SVM sucht eine Trennlinie mit möglichst grossem Margin zwischen den Klassen.
Der Parameter $C$ steuert die Fehlertoleranz: grosses $C$ für strenge Trennung, kleines $C$ für mehr Toleranz.
Mit Kernels lassen sich auch nichtlinear trennbare Daten behandeln.
\end{cheatSheetBox}

\columnbreak\textcolor{\cheatSheetColor}{\subsection*{Assoziation und Filterung}}

\begin{cheatSheetBox}[title=Apriori]
Apriori sucht häufige Itemsets und daraus Assoziationsregeln.

\[
\text{Support}(X\Rightarrow Y)=P(X\cup Y)
\qquad
\text{Confidence}(X\Rightarrow Y)=\frac{P(X\cup Y)}{P(X)}
\qquad
\text{Lift}(X\Rightarrow Y)=\frac{P(X\cup Y)}{P(X)P(Y)}
\]

Hohe Confidence allein genügt nicht; Support und Lift zeigen, ob die Regel wirklich sinnvoll ist.
\end{cheatSheetBox}

\begin{cheatSheetBox}[title=Naive Bayes]
Naive Bayes kombiniert Bayes' Theorem mit der Annahme bedingter Unabhängigkeit der Merkmale.

\[
P(C\mid x)=\frac{P(x\mid C)P(C)}{P(x)}
\qquad
P(x_1,\ldots,x_n\mid C)\approx\prod_{i=1}^{n}P(x_i\mid C)
\]

Laplace-Glättung verhindert Nullwahrscheinlichkeiten und ist besonders wichtig bei Textdaten.
\end{cheatSheetBox}

\columnbreak\textcolor{\cheatSheetColor}{\subsection*{Regression und Evaluation}}

\begin{cheatSheetBox}[title=Regression]
Regression sagt einen kontinuierlichen Zahlenwert vorher.

\[
y=a_0+a_1x_1+a_2x_2+\cdots+a_nx_n
\]

\[
\text{MAE}=\frac{1}{n}\sum_{i=1}^{n}|y_i-\hat{y}_i|
\qquad
\text{RMSE}=\sqrt{\frac{1}{n}\sum_{i=1}^{n}(y_i-\hat{y}_i)^2}
\]

RMSE bestraft grosse Fehler stärker als MAE.
\end{cheatSheetBox}

\begin{cheatSheetBox}[title=Datenvorbereitung]
\begin{itemize}
    \item Train-Test-Split: typischerweise 80/20.
    \item Features sind Eingaben, das Label ist die Zielvariable.
    \item Datenleckage verfälscht die Bewertung.
    \item Klassifikation: Accuracy, Precision, Recall.
    \item Regression: MAE, RMSE.
\end{itemize}
\end{cheatSheetBox}

\begin{cheatSheetBox}[title=Merksätze]
\begin{itemize}
    \item Erst Daten sauber vorbereiten, dann trainieren.
    \item Ein Modell ohne gute Evaluation ist wertlos.
    \item Interpretierbarkeit und Güte müssen zusammen betrachtet werden.
\end{itemize}
\end{cheatSheetBox}

\end{multicols*}
\end{landscape}
\clearpage
\restoregeometry